\section{Semantic Model}
\label{sec:semantic-model}

The semantic model defines the data types that flow through the system.

\subsection{SemanticFact}
\label{sec:semantic-model:fact}

The core data type is \texttt{SemanticFact}, an immutable record with the
following fields:

\begin{description}
    \item[id] A globally unique identifier (UUID).
    \item[subject] The entity identifier that is the subject of the fact.
    \item[relation] The relation type identifier.
    \item[objects] An ordered tuple of values (entity IDs or literals).
    \item[context] The context in which this fact is valid.
    \item[confidence] A floating-point value in $[0, 1]$ indicating
          certainty.
    \item[metadata] An optional map of key-value pairs for extensibility.
\end{description}

\subsection{Entity and Relation}
\label{sec:semantic-model:entity}

Entities and relations are typed identifiers. Each entity has a semantic
type (e.g., \texttt{Person}, \texttt{Organization}, \texttt{Event}) and
may carry attributes. Each relation has an arity and typed slots.

\subsection{Context}
\label{sec:semantic-model:context}

A context is a dot-separated path string representing a scope of validity.
Contexts form a finite poset under prefix inclusion (Definition~\ref{def:context}).
The root context ``world'' is the least specific; sub-contexts refine the scope.

\subsection{Invariants}
\label{sec:semantic-model:invariants}

Every \texttt{SemanticFact} satisfies:
\begin{enumerate}
    \item Immutability: once created, fields do not change.
    \item Identifier uniqueness: no two facts share an ID.
    \item Confidence bounds: $0.0 \leq \text{confidence} \leq 1.0$.
    \item Referential integrity: subject and referenced objects exist.
\end{enumerate}
